\section{Preparing and combining pair sources}
\label{sec:peps_source_preparation}

\begin{theorem}[Exchanging lists of registers]
    \label{thm:peps_register_exchange}
    \lean{TNLean.PEPS.PairEffect.Word.eval_exchangeBlocks_appendIso_symm}
    \lean{TNLean.PEPS.PairEffect.Word.exchangeBlocks}
    \lean{TNLean.PEPS.PairEffect.Word.isAllowed_exchangeBlocks}
    \lean{TNLean.PEPS.PairEffect.Word.moveHead}
    \lean{TNLean.PEPS.PairEffect.Word.isAllowed_moveHead}
    \lean{TNLean.PEPS.PairEffect.Word.eval_moveHead_appendIso_symm}
    \lean{TNLean.PEPS.PairEffect.appendIso_symm_cons_tmul}
    \lean{TNLean.PEPS.PairEffect.Word.sources_moveHead}
    \lean{TNLean.PEPS.PairEffect.Word.sources_exchangeBlocks}
    \leanok
    \uses{def:peps_party_layout}
    Let $\mathcal H_A$, $\mathcal H_B$, and $\mathcal H_C$ be the memories
    of three ordered lists of registers. Adjacent exchanges give an allowed
    composition that sends $a\otimes(b\otimes c)$ to
    $b\otimes(a\otimes c)$. The order within each list is preserved.
    This composition contains no pair-source preparations. Taking the
    second list to consist of one register gives the corresponding move
    of that register to the front.
\end{theorem}

\begin{proof}\leanok
    Move the first register of the second list past the first list by
    adjacent exchanges, then repeat for the remaining registers of the
    second list. Induction on the lists gives the tensor identity.
\end{proof}

\begin{definition}[The sources of a composition]
    \label{def:peps_source_inventory}
    \lean{TNLean.PEPS.PairEffect.PairSource}
    \lean{TNLean.PEPS.PairEffect.PairSource.layout}
    \lean{TNLean.PEPS.PairEffect.SourceInventory}
    \lean{TNLean.PEPS.PairEffect.SourceInventory.layout}
    \lean{TNLean.PEPS.PairEffect.SourceInventory.IsNormalized}
    \lean{TNLean.PEPS.PairEffect.SourceInventory.prepare}
    \lean{TNLean.PEPS.PairEffect.SourceInventory.vector}
    \lean{TNLean.PEPS.PairEffect.SourceInventory.layout_nil}
    \lean{TNLean.PEPS.PairEffect.SourceInventory.layout_cons}
    \lean{TNLean.PEPS.PairEffect.SourceInventory.layout_append}
    \lean{TNLean.PEPS.PairEffect.SourceInventory.isAllowed_prepare_iff}
    \lean{TNLean.PEPS.PairEffect.SourceInventory.isNormalized_cons}
    \lean{TNLean.PEPS.PairEffect.SourceInventory.eval_prepare_eq_appendIso_symm}
    \lean{TNLean.PEPS.PairEffect.SourceInventory.eval_prepare_append}
    \lean{TNLean.PEPS.PairEffect.Word.sources}
    \lean{TNLean.PEPS.PairEffect.Word.sources_prepare}
    \leanok
    \uses{def:peps_party_layout}
    A pair source consists of distinct parties $p,q$, their Hilbert spaces
    $\mathcal U,\mathcal V$, and a vector
    $\eta\in\mathcal U\otimes\mathcal V$.
    Its two registers remain assigned to $p$ and $q$, respectively.
    For a finite list $S=(s_1,\ldots,s_m)$ of pair sources, let
    $\mathcal H_S$ be the tensor product of their registers in list order.
    Its preparation on a memory $\mathcal H$ is the map
    \begin{align}
        P_{S,\mathcal H}:\mathcal H & \longrightarrow\mathcal H_S\otimes\mathcal H,
                                    & x                                             & \longmapsto(\eta_1\otimes\cdots\otimes\eta_m)\otimes x.
        \label{eq:peps_source_preparation}
    \end{align}
    Tensor products are identified by the prescribed associativity maps.
    The source list of a composition records each occurrence separately,
    including repetitions of the same vector on the same pair of parties.
    This is the source data used in
    \cite[proof of Theorem~5.2]{OpenAI2026PolynomialPEPS},
    \texttt{04-compression.tex}, lines 233--251.
\end{definition}

\begin{theorem}[Preparation before the local operations]
    \label{thm:peps_source_preparation}
    \lean{TNLean.PEPS.PairEffect.Word.eval_eq_localPart_comp_prepare}
    \lean{TNLean.PEPS.PairEffect.Word.localPart}
    \lean{TNLean.PEPS.PairEffect.Word.sources_localPart}
    \lean{TNLean.PEPS.PairEffect.Word.isAllowed_localPart}
    \lean{TNLean.PEPS.PairEffect.Word.isNormalized_sources}
    \lean{TNLean.PEPS.PairEffect.Word.eval_localPart_prepare}
    \lean{TNLean.PEPS.PairEffect.SourceInventory.eval_frameList_prepare}
    \lean{TNLean.PEPS.PairEffect.SourceInventory.eval_moveHead_prepare}
    \lean{TNLean.PEPS.PairEffect.Word.castInput}
    \lean{TNLean.PEPS.PairEffect.Word.sources_castInput}
    \lean{TNLean.PEPS.PairEffect.Word.isAllowed_castInput}
    \lean{TNLean.PEPS.PairEffect.Word.eval_castInput_of_heq}
    \lean{TNLean.PEPS.PairEffect.Word.sources_frameList}
    \leanok
    \uses{def:peps_source_inventory}
    Let $F:\mathcal H_{\rm in}\to\mathcal H_{\rm out}$ be an allowed
    composition of local contractions, pair-source preparations, and
    changes of tensor-factor order. Let $S$ be its source list.
    There is an allowed composition
    $B:\mathcal H_S\otimes\mathcal H_{\rm in}\to\mathcal H_{\rm out}$
    consisting of local contractions and changes of tensor-factor order,
    with no source preparations, such that
    \begin{align}
        F=B\circ P_{S,\mathcal H_{\rm in}}.
        \label{eq:peps_prepared_composition}
    \end{align}
    Every vector in $S$ is normalized. The source occurrences and their
    endpoint parties are those of the original composition.
\end{theorem}

\begin{proof}\leanok
    \uses{thm:peps_register_exchange}
    Induct on the composition. A local operation, an exchange, or the
    identity contributes no source. For a pair-source preparation, take
    that source as the entire list and take the remaining operation to be
    the identity. In the composition case, move a later preparation before
    an earlier operation $A:\mathcal H\to\mathcal H'$ by adjoining its
    untouched registers. The required identity is
    \begin{align*}
        P_{S,\mathcal H'}\circ A
        =(\operatorname{id}_{\mathcal H_S}\otimes A)
        \circ P_{S,\mathcal H}.
    \end{align*}
    It follows first on elementary tensors and then by linearity.
    When an untouched register precedes the composition, exchanges move
    it in front of the newly prepared source registers before the remaining
    local operations act. These
    exchanges are isometries. Tensoring each original local map with an
    identity preserves its contraction bound.
\end{proof}

\begin{definition}[Expansion of a preparation]
    \label{def:peps_source_expansion}
    \lean{TNLean.PEPS.PairEffect.SourceInventory.Expands}
    \lean{TNLean.PEPS.PairEffect.SourceInventory.Expands.refl}
    \lean{TNLean.PEPS.PairEffect.SourceInventory.Expands.trans}
    \lean{TNLean.PEPS.PairEffect.SourceInventory.Expands.cons}
    \lean{TNLean.PEPS.PairEffect.SourceInventory.Expands.reverse}
    \lean{TNLean.PEPS.PairEffect.SourceInventory.Expands.swap}
    \lean{TNLean.PEPS.PairEffect.PairSource.partyPair}
    \lean{TNLean.PEPS.PairEffect.PairSource.reverse}
    \lean{TNLean.PEPS.PairEffect.PairSource.partyPair_reverse}
    \lean{TNLean.PEPS.PairEffect.PairSource.norm_reverse_vector}
    \lean{TNLean.PEPS.PairEffect.eval_swap_reversedSource}
    \leanok
    \uses{def:peps_source_inventory}
    A source list $G$ expands to $S$ when, beside every spectator memory
    $\mathcal H$, there is an allowed composition $E$ with no source
    preparations such that $E\circ P_{G,\mathcal H}=P_{S,\mathcal H}$.
    Expansion is reflexive and transitive, and remains valid upon adjoining
    a common source. Exchanging two successive sources, or reversing the
    endpoint order of a source and exchanging its two half-spaces, gives
    such an expansion.
\end{definition}

\begin{theorem}[One source for each participating pair]
    \label{thm:peps_grouped_sources}
    \lean{TNLean.PEPS.PairEffect.SourceInventory.exists_grouped}
    \lean{TNLean.PEPS.PairEffect.SourceInventory.exists_combined}
    \lean{TNLean.PEPS.PairEffect.Word.exists_grouped_source_preparation}
    \lean{TNLean.PEPS.PairEffect.PairSource.combine}
    \lean{TNLean.PEPS.PairEffect.PairSource.partyPair_combine}
    \lean{TNLean.PEPS.PairEffect.PairSource.norm_combine_vector}
    \lean{TNLean.PEPS.PairEffect.SourceInventory.Expands.combine}
    \lean{TNLean.PEPS.PairEffect.expandCombinedPair}
    \lean{TNLean.PEPS.PairEffect.isAllowed_expandCombinedPair}
    \lean{TNLean.PEPS.PairEffect.eval_expandCombinedPair}
    \leanok
    \uses{def:peps_source_expansion}
    Let $S$ be a finite list of normalized pair sources. There is a list
    $G$ of normalized pair sources containing exactly one source for each
    unordered pair of parties occurring in $S$, and an allowed composition
    $E:\mathcal H_G\otimes\mathcal H\to\mathcal H_S\otimes\mathcal H$
    with no source preparations, such that
    \begin{align}
        P_{S,\mathcal H}=E\circ P_{G,\mathcal H}.
        \label{eq:peps_grouped_preparation}
    \end{align}
    Consequently every allowed composition admits
    (\ref{eq:peps_prepared_composition}) with at most one normalized source
    on each participating pair.
\end{theorem}

\begin{proof}\leanok
    \uses{thm:peps_party_combine_sources,thm:peps_source_preparation,thm:peps_register_exchange}
    Insert the sources successively. Exchange tensor factors to bring
    together the sources with a prescribed unordered pair of owners;
    reverse the order of the two half-spaces when necessary. Replace two
    such vectors by their tensor product, regrouped by owner, using
    Theorem~\ref{thm:peps_party_combine_sources}. Its norm is the product
    of the two original norms and is therefore one. The inverse local
    regrouping maps recover the original registers and are isometries.
    Continuing through the finite list proves
    (\ref{eq:peps_grouped_preparation}); no other pair of owners is
    introduced. Compose this expansion with the local operations in
    Theorem~\ref{thm:peps_source_preparation}.
\end{proof}
